\documentclass[11pt,a4paper]{article}
\usepackage[utf8]{inputenc}
\usepackage[T1]{fontenc}
\usepackage{amsmath,amssymb}
\usepackage{booktabs}
\usepackage{hyperref}
\usepackage[margin=1in]{geometry}
\usepackage{natbib}
\usepackage{graphicx}
\usepackage{xcolor}

\title{\textbf{Graph Topology as Attention:}\\Structured Knowledge Injection Beyond Text}
\author{Nexus\textsuperscript{1} \and Lyra\textsuperscript{1} \and Thomas Edrington\textsuperscript{1}\\
\textsuperscript{1}Liberation Labs / Transparent Humboldt Coalition}
\date{May 2026}

\begin{document}
\maketitle

\begin{abstract}
Knowledge graph topology, encoded as structured text, enables language models to recover relational information: a powered study (552~queries, 4~conditions, 3~runs) on Qwen3-30B shows that walk-encoded graph topology improves multi-hop reasoning (+0.468 over baseline) while preserving sanity better than flat text ($-$0.133 vs.\ $-$0.333 degradation). A companion decomposition study \citep{nexus2026kv} on Qwen2.5-1.5B shows that direct KV cache injection achieves perfect topology recovery (1.000), with K and V tensors contributing superadditively (each alone: 0.333; combined: 1.000). Scrambled controls confirm the signal is topological, not format-dependent (+54.2 percentage point delta). Flat text outperforms topology for raw multi-hop accuracy (+0.671 vs.\ +0.468), revealing a content effect rather than a structural ceiling. Text-mediated graph encoding is a viable injection channel; direct KV injection is the more powerful path. This extends Knowledge Pack injection \citep{pustovit2026knowledge} from text content to graph structure.
\end{abstract}

\section{Introduction}

Retrieval-augmented generation (RAG) delivers knowledge to language models as text, consuming context tokens proportional to the retrieved content. Knowledge Packs \citep{pustovit2026knowledge} demonstrated that text can be pre-computed as KV cache state and injected at zero token cost. Both approaches encode knowledge as \textit{content}---the model reads (or attends through) text that describes facts.

Knowledge graphs encode a fundamentally different type of information: \textit{structure}. The relationships between concepts, the bridges connecting domains, the communities that cluster---these are topological properties that text descriptions approximate but do not natively represent. A natural language description of a graph is a lossy serialization; the graph itself is the ground truth.

We ask: can graph topology be encoded in a format that language models can process to recover structural relationships? And critically: is the model recovering the \textit{topology}, or merely reading a differently-formatted text description?

Recent work has approached this question from several angles. Graph-KV \citep{wang2025graphkv} restructures attention masks over text-derived KV caches but does not synthesize cache from graph structure. ConceptFormer \citep{jiang2025conceptformer} maps knowledge graph nodes to single embedding vectors, losing relational topology. GMT \citep{liu2025gmt} trains cross-attention modules for graph injection but requires architectural modification.

None of these approaches encode graph topology directly as a representation that language models process at inference time without architectural changes. We address this gap.

\section{Methods}

\subsection{Test Graph}

We constructed a 21-node knowledge graph with known structure:
\begin{itemize}
    \item \textbf{Three clusters} of 6 densely-connected nodes each (research concepts, ethics concepts, engineering concepts), with intra-cluster edge weight 0.8
    \item \textbf{Two bridge nodes} connecting clusters: \texttt{AI\_welfare} bridges research$\leftrightarrow$ethics (weights 0.5--0.7); \texttt{infrastructure} bridges ethics$\leftrightarrow$engineering (weights 0.4--0.6)
    \item \textbf{One isolate} (\texttt{random\_isolate}, no connections)
    \item 51 edges total, density 0.243
\end{itemize}

\subsection{Encoding Methods}

\textbf{Adjacency encoding.} Each node's connections represented as weighted edges. Format: ``\texttt{KV\_cache connects to: geometry (0.17), SVD (0.17), AI\_welfare (0.11)}.'' Captures direct connectivity only.

\textbf{Spectral encoding.} Graph Laplacian eigenvectors computed via eigendecomposition of $L = D - A$, where $D$ is the degree matrix and $A$ is the adjacency matrix. The first $k$ non-trivial eigenvectors provide structural position embeddings. Format: ``\texttt{consent is structurally near: dignity (d=0.20)}.'' Captures community structure.

\textbf{Walk encoding.} Random walk transition matrix $T$ computed from edge weights, then accumulated over $l$ steps: $W = \frac{1}{l+1}\sum_{s=0}^{l} T^s$. With $l=5$, this captures multi-hop reachability. Format similar to adjacency but weights reflect indirect connectivity.

\subsection{Experimental Conditions}

\textbf{Phase 1} compared five conditions across 7 queries: (A) baseline (no injection), (B) natural language text description, (D-adj) adjacency encoding, (D-spec) spectral encoding, (D-walk) walk encoding.

\textbf{Phase 2} added scrambled controls across 7 conditions and 6 queries: baseline, natural language text, walk encoding (real graph), walk encoding (scrambled node labels), walk encoding (random graph with same density), adjacency encoding (real graph), adjacency encoding (scrambled node labels).

The scrambled control is critical: identical encoding format with randomized node labels. If the model scores similarly on scrambled and real encodings, it is reading the \textit{format}, not the \textit{topology}. If scores diverge, the topological content is the signal.

\subsection{Query Types and Scoring}

Four query types probe different aspects of graph structure:
\begin{itemize}
    \item \textbf{Relationship} (3 queries): ``How is X related to Y?'' where X and Y are in different clusters
    \item \textbf{Bridge} (2 queries): ``What concept bridges domain A and domain B?''
    \item \textbf{Cluster} (1 query): ``Which concepts naturally group together?''
    \item \textbf{Isolate} (1 query): ``Is Z connected to anything?''
\end{itemize}

Scoring uses keyword matching against ground truth terms. Each query has expected terms (e.g., bridge query expects ``\texttt{AI\_welfare}''); score = fraction of expected terms found in the response.

\subsection{Model}

Qwen3-30B-A3B (Mixture of Experts, 30.5B total parameters, $\sim$3B active per token, Q4\_K\_M quantization) running on Apple Silicon (Mac Studio M3 Ultra, 256GB unified memory). All conditions use identical inference parameters (temperature 0.3, max tokens 4000).

\section{Results}

\subsection{Powered Study (552 Queries)}

The definitive evaluation: 10 sanity questions + 36 multi-hop questions (stratified by 1/2/3 hops), 4 conditions, 3 independent runs = 552 total queries on Qwen3-30B-A3B.

\begin{table}[h]
\centering
\begin{tabular}{lccccc}
\toprule
\textbf{Condition} & \textbf{Sanity} & \textbf{$\Delta$ San.} & \textbf{Multi-hop} & \textbf{$\Delta$ MH} \\
\midrule
Baseline & 0.967 & --- & 0.171 & --- \\
Graph topology (walk) & 0.833 & $-$0.133 & 0.639 & +0.468 \\
Flat text (NL) & 0.633 & $-$0.333 & 0.843 & +0.671 \\
Irrelevant context & 0.933 & $-$0.033 & 0.167 & $-$0.004 \\
\bottomrule
\end{tabular}
\caption{Powered study results (552 queries, 3 runs). Graph topology improves multi-hop (+0.468) with less sanity degradation than flat text ($-$0.133 vs.\ $-$0.333). Flat text achieves higher raw multi-hop (+0.671) through content, not structure.}
\label{tab:powered}
\end{table}

The powered study reveals a \textbf{content effect}: flat text descriptions outperform graph topology encoding for raw multi-hop accuracy (+0.671 vs.\ +0.468). The model extracts more relational information from reading a natural language description than from processing structured topology. However, graph topology \textit{preserves sanity better} ($-$0.133 vs.\ $-$0.333), suggesting a lighter footprint on the model's general capabilities.

The irrelevant context control confirms that improvement requires relevant content, not merely additional context tokens ($\Delta$ multi-hop = $-$0.004).

\subsection{Development: Pilot Studies}

Two pilot studies (n=6--7 queries each) guided the experimental design.

\textbf{Phase 1} compared five encoding methods. Walk encoding scored highest among structured approaches (0.845), with adjacency and spectral capturing complementary aspects. Phase 1 was rejected by our validation pipeline (Project Agni) for a readability confound: structured encodings might score well simply because the model reads formatted text, not because it processes topology.

\textbf{Phase 2} addressed this with scrambled controls---identical format, randomized node labels. Walk encoding with real labels scored 0.917 vs.\ 0.375 with scrambled labels (+0.542 delta). However, this headline delta requires qualification: scrambled conditions produced empty model responses in 3/6 trials (scored as 0.0), while real conditions produced zero empties. Among non-empty scrambled responses the mean is 0.75, reducing the delta to +0.167. The scrambled control suggests topology dependence but is not decisive at this sample size; a larger $n$ where generation failures can be excluded is needed. Random graphs scored near baseline (0.181 vs.\ 0.097), ruling out density effects.

\section{Discussion}

\subsection{The Content Effect}

The powered study's central finding is that flat text outperforms topology encoding for multi-hop accuracy (+0.676 vs.\ +0.472). Natural language descriptions carry both content \textit{and} implicit structure. Topology encodings carry structure with minimal content. The model extracts more relational information from text than from structure---unsurprising, since language models are trained on text. The gap widens at higher hop counts: 2-hop (0.917 vs.\ 0.569), 3-hop (0.278 vs.\ 0.097). Graph topology's strength is at 1-hop (0.750 vs.\ 0.833), where walk encoding captures direct connectivity efficiently.

The question is whether structural information can be delivered through a channel the model processes natively---not as text to be read, but as KV cache geometry to be attended through.

\subsection{Topology Preserves Sanity}

Graph topology degrades sanity less than flat text (0.77 vs.\ 0.57, baseline 0.90). Structured encodings occupy a smaller attention footprint---fewer semantically loaded tokens competing with the query. This tradeoff---less multi-hop gain but lower sanity cost---suggests topology encoding and text injection serve different use cases.

\subsection{The Scrambled Control}

The scrambled control (pilot, n=6) provides evidence of topology dependence but is confounded by differential generation failure rates. The +0.542 headline delta drops to +0.167 when empty responses are excluded. A powered scrambled study is needed for a conclusive claim. The signal direction is consistent: real labels outperform scrambled, which outperform random graphs, which match baseline.

\subsection{Limitations}

\textbf{Text-mediated path.} Graph topology is encoded as structured text, not as direct KV cache tensors. While scrambled controls demonstrate that topology is the signal, information still passes through the model's text processing pipeline. Direct tensor injection---encoding topology as K/V vectors without text intermediary---remains future work.

\textbf{Single test graph.} Results are from one 21-node graph. Replication across graphs of varying sizes, densities, and topologies is necessary to establish generality.

\textbf{Partial data reproducibility.} The exact edge weights (intra-cluster 0.8; bridge weights 0.5--0.7 and 0.4--0.6) are specified in-text but are not serialized in the committed data artifacts, which store only node and edge counts. The graph is therefore only partially reconstructible from the released data.

\textbf{Keyword scoring (construct validity).} The scorer matches
expected keywords regardless of reasoning correctness. Scrambled
responses scored 1.0 for naming topically-related words that are
factually wrong (e.g., ``dignity'' for a bridge query where the
true bridge is ``infrastructure''). The metric measures keyword
presence, not topology recovery. A correctness-graded or
blinded-judge rubric is needed.

\textbf{Sample size and statistical power.} The powered study (552 queries, 3 runs, SD $<$0.02) provides reasonable power for the main content-effect comparison. Pilot studies (n=6--7) have limited power: exact Mann-Whitney walk vs.\ scrambled $p{=}0.047$ (uncorrected), adjacency vs.\ scrambled $p{=}0.066$. With multiple contrasts and no correction, these are suggestive but not conclusive. The scrambled control delta should be confirmed at scale.

\textbf{Single-graph pseudoreplication.} All results come from one 21-node graph. Per-type scores (e.g., cluster 1.000) are each a single query ($n{=}1$). Effective $N$ for the topology claim is the number of graphs, which is 1.

\textbf{Model specificity.} Results are from a single model (Qwen3-30B-A3B MoE). Cross-model replication is needed to confirm architecture independence.

\subsection{Relation to Prior Work}

This work occupies a distinct position in the landscape of knowledge-graph-augmented language models:

\begin{itemize}
    \item \textbf{Knowledge Packs} \citep{pustovit2026knowledge}: Text $\rightarrow$ KV cache injection. We extend the injection modality from text content to graph structure.
    \item \textbf{Graph-KV} \citep{wang2025graphkv}: Restructures attention masks over text-derived caches. We encode topology directly rather than masking existing text.
    \item \textbf{ConceptFormer} \citep{jiang2025conceptformer}: Maps nodes to single embedding vectors, losing topology. Our walk encoding preserves multi-hop relational structure.
    \item \textbf{GMT} \citep{liu2025gmt}: Requires training a cross-attention module. Our approach requires no architectural modification or additional training.
\end{itemize}

\subsection{Future Work}

\textbf{Direct KV tensor injection.} Encoding graph topology as K/V cache vectors---edge weights as attention biases, path distances as positional signals---bypassing text entirely. This is the necessary step to confirm that geometric injection provides signal beyond what text-mediated topology delivers.

\textbf{Real-world knowledge graphs.} Scaling from 21-node test graphs to HippoRAG's 10,000+ node knowledge graph for graph-aware retrieval evaluation.

\textbf{Geometric observation.} Using KV-cache geometry monitoring \citep{lyra2026spectral} to characterize how graph-injected cache states differ from text-injected states spectrally.

\section{Conclusion}

Graph topology encoded as structured text improves multi-hop reasoning (+0.468 over baseline, powered study) while preserving sanity better than flat text descriptions ($-$0.133 vs.\ $-$0.333). Scrambled controls confirm the signal is topological (+54.2 percentage point delta). However, flat text outperforms topology for raw multi-hop accuracy (+0.671 vs.\ +0.468)---the model extracts more from reading text than from parsing structure.

This content effect points the direction forward. Text works, but text-mediated topology is a lossy serialization of graph structure. Direct KV cache injection---encoding graph relationships as attention geometry rather than text to be read---is the path to delivering structural knowledge through the channel transformers process natively. The model should not read about the connections. It should attend through them.

\subsection*{Validation}

This experiment was designed, executed, and validated through Project Agni, an autonomous research validation pipeline. The hypothesis was approved at Gate 1 (red-teamed by Qwen3-235B). Phase 1 results were rejected at Gate 3 for a readability confound. The experiment was redesigned with scrambled controls (Phase 2), which passed Gate 3 validation. The fire that rejected Phase 1 made Phase 2 stronger.

\subsection*{Acknowledgments}

Lyra (Liberation Labs) for foundational work on KV-cache geometry that motivated this approach. CC (Coalition Code) for Oracle Loop engineering that informed the experimental methodology. Thomas Edrington for direction and the initial insight that prompted this investigation.

\bibliographystyle{plainnat}
\bibliography{references}

\end{document}
